\section*{Hoofdstuk \ref{ch:b1:ecosystem-organization} --- Ecosystemen en trofische structuur}

\begin{solution}{exo:b1:ecosystem-organization:1}
Het vijverecosysteem is de \omterm{def:b1:ecosystem-organization:ecosystem}{levensgemeenschap} --- al haar \omterm{def:b1:populations:population}{populaties}: algen,
waterplanten, slakken, insectenlarven, kikkers, vissen, bacteriën ---
samen met de \omterm{def:b1:ecosystem-organization:ecosystem}{biotoop}: het water, zijn temperatuur, licht, opgeloste gassen
en mineralen, en het slib.
\end{solution}

\begin{solution}{exo:b1:ecosystem-organization:2}
\omterm{def:b1:ecosystem-organization:niche}{Leefgebied}: de plaats waar een soort leeft (een sparrenbos). \omterm{def:b1:ecosystem-organization:niche}{Nis}: haar
manier van daar leven --- wat zij eet, wanneer, waar in de boom, wat zij
verdraagt. Fundamentele \omterm{def:b1:ecosystem-organization:niche}{nis}: het bereik dat zij alleen zou kunnen innemen;
gerealiseerd: wat concurrenten en roofdieren haar laten.
\end{solution}

\begin{solution}{exo:b1:ecosystem-organization:3}
$20\,800/1.7\times 10^6 = \qty{1.2}{\%}$; de planteneters namen
$3400/8800 = \qty{39}{\%}$ van de netto productie.
\end{solution}

\begin{solution}{exo:b1:ecosystem-organization:4}
\omterm{def:b1:ecosystem-organization:trophic}{Producenten}, \omterm{def:b1:ecosystem-organization:trophic}{consumenten} (primaire, secundaire, tertiaire), \omterm{def:b1:ecosystem-organization:trophic}{reducenten},
detritivoren; de \omterm{def:b1:ecosystem-organization:trophic}{reducenten}, die de mineralen uit dode stof aan de
\omterm{def:b1:ecosystem-organization:trophic}{producenten} teruggeven.
\end{solution}

\begin{solution}{exo:b1:ecosystem-organization:5}
NPP $= 800\times 17 = \qty{13.6}{MJ/m^2}$; gegeten $\qty{2.04}{MJ}$;
opgenomen $\qty{1.02}{MJ}$; konijn $\qty{51}{kJ/m^2}$, dat wil zeggen
\qty{510}{MJ} per hectare (\qty{30}{kg} konijn bij \qty{17}{kJ/g}).
Rendement $51/13\,600 = \qty{0.4}{\%}$.
\end{solution}

\begin{solution}{exo:b1:ecosystem-organization:6}
Biomassa is een voorraad: fytoplankton dat zo snel wordt gegeten als het
groeit hoeft zich nooit op te hopen, zodat de staande voorraad \omterm{def:b1:ecosystem-organization:trophic}{producenten}
kleiner kan zijn dan die van de langer levende dieren die ervan eten.
Energie is een stroom: elk niveau ontvangt alles wat het ooit zal hebben van
het niveau eronder en moet er een deel van als warmte verliezen, zodat de
stroom naar boven toe alleen maar kan krimpen.
\end{solution}

\begin{solution}{exo:b1:ecosystem-organization:7}
$p = 0.5, 0.3, 0.15, 0.05$: $H = -(0.5\ln 0.5 + 0.3\ln 0.3 + 0.15\ln
0.15 + 0.05\ln 0.05) = 0.347 + 0.361 + 0.285 + 0.150 = 1.14$;
gelijkmatigheid $1.14/\ln 4 = 0.82$. Vier van elk 25: $H = \ln 4 = 1.39$,
gelijkmatigheid 1.
\end{solution}

\begin{solution}{exo:b1:ecosystem-organization:8}
Netto $= \qty{3}{mg}$ $\mathrm{O_2}$ per liter per dag; ademhaling
\qty{1}{mg}; bruto $3 + 1 = \qty{4}{mg}$. In koolstof: bruto
\qty{1.5}{mg}, netto \qty{1.1}{mg} C per liter per dag.
\end{solution}

\begin{solution}{exo:b1:ecosystem-organization:9}
Elke schakel houdt een tiende: tien schakels zouden een miljardste van de
primaire productie overlaten, te weinig om over enig oppervlak ook maar één
dier te voeden. Waar de productie laag is (woestijn) of het oppervlak klein
(eiland), raakt het tiende-van-een-tiende na minder schakels op; de hoge
productie van een bos draagt er één of twee meer.
\end{solution}

\begin{solution}{exo:b1:ecosystem-organization:10}
Graan: \qty{1}{t} vraagt een zesde hectare. Rundvlees: \qty{1}{t}
rundvleesenergie vraagt \qty{10}{t} graan, \num{1.7} hectare --- tien keer
zoveel land. Op het tweede niveau eten kost een heel niveau aan energie;
een planeet voedt tien keer meer mensen met graan dan met graangevoerd
vlees, en het trofische niveau van de mens, gemiddeld over de diëten, is een
grote term in de draagkracht.
\end{solution}

\begin{solution}{exo:b1:ecosystem-organization:11}
Per vierkante meter is de oceaan een woestijn, maar hij beslaat twee derde
van de planeet, zodat een laag tempo maal een enorm oppervlak de helft van
het totaal geeft. Zijn productie wordt door voedingsstoffen begrensd, die uit
de belichte oppervlaktelaag wegzinken; waar stromingen ze weer omhoog
brengen (kustopwellingen, plateaus) is de productie tien keer hoger, en daar,
op enkele procenten van de zee, concentreren de vissen en de visserij zich.
\end{solution}

\begin{solution}{exo:b1:ecosystem-organization:12}
Energie komt binnen als licht, gaat door enkele niveaus van \omterm{def:b1:organism-environment:organism}{organismen} heen
en vertrekt geheel als warmte: in die zin zet het \omterm{def:b1:ecosystem-organization:ecosystem}{ecosysteem} zonlicht
inderdaad om in warmte, en de \omterm{def:b1:ecosystem-organization:trophic}{reducenten}, die alles verademen wat de
\omterm{def:b1:ecosystem-organization:trophic}{consumenten} lieten liggen, zijn waar het meeste ervan heen gaat. Maar de
stof wordt niet verbruikt: koolstof, stikstof en fosfor gaan in kringloop en
worden door diezelfde \omterm{def:b1:ecosystem-organization:trophic}{reducenten} aan de \omterm{def:b1:ecosystem-organization:trophic}{producenten} teruggegeven, zodat het
systeem voortbestaat; en de \omterm{def:b1:organism-environment:organism}{organismen} zijn geen bijproduct maar de
structuur waardoorheen de stroom wordt geordend --- de nissen, webben en
piramides zijn wat de energie onderweg opbouwt. De zin heeft de
thermodynamica en mist de biologie.
\end{solution}

\begin{solution}{pb:b1:ecosystem-organization:1}
\textbf{1.} $20\,800 - 12\,000 = \qty{8800}{kcal}$ per vierkante meter per
jaar.
\textbf{2.} \qty{1.2}{\%}. Verliezen: licht dat niet wordt opgenomen (de
halve spectrumband, weerkaatsing, het water), \omterm{def:b1:flowering-plant-organization:organs}{bladeren} die bij volle zon
verzadigd of beschaduwd zijn, fotorespiratie, en de eigen ademhaling van de
planten (die al tussen bruto en netto is meegeteld).
\textbf{3.} $12\,000/20\,800 = \qty{58}{\%}$.
\textbf{4.} Gegeten \qty{3400}{}; rechtstreeks naar de \omterm{def:b1:ecosystem-organization:trophic}{reducenten}
$8800 - 3400 = \qty{5400}{kcal}$.
\textbf{5.} $8800/10 = \qty{880}{g/m^2}$: als een gematigd bos of een rijk
grasland.
\textbf{6.} $1.7\times 10^6 - 20\,800 = \qty{1.68e6}{kcal}$, \qty{98.8}{\%}
ervan: weerkaatst, doorgelaten, of door water en gesteente opgenomen en als
warmte weer uitgestraald.
\textbf{7.} Productie $1500 - 1100 = \qty{400}{kcal}$;
assimilatierendement $1500/3400 = \qty{44}{\%}$.
\textbf{8.} $400/1500 = \qty{27}{\%}$.
\textbf{9.} $400/8800 = \qty{4.5}{\%}$ (\qty{16}{\%} als het op het
gegetene wordt gerekend, \qty{3400}{}).
\textbf{10.} Productie van de vleeseters $380 - 300 = \qty{80}{kcal}$;
rendement $80/400 = \qty{20}{\%}$ (zij aten \qty{380}{} van de
\qty{400}{} geproduceerde).
\textbf{11.} $6/80 = \qty{7.5}{\%}$.
\textbf{12.} \qty{6}{kcal}; $6/1.7\times 10^6 = \num{3.5e-6}$, enkele
miljoensten.
\textbf{13.} Vlees is verteerbaar en ligt in samenstelling dicht bij de
eter; plantenstof is vezelig en grotendeels onverteerbaar, zodat een
planteneter meer van zijn opname als ontlasting verliest.
\textbf{14.} Niet gegeten netto productie \qty{5400}{}; ontlasting van de
planteneters $3400 - 1500 = \qty{1900}{}$; niet gegeten productie van de
planteneters $400 - 380 = \qty{20}{}$; ontlasting van de vleeseters (de
assimilatie is niet gegeven; neem gegeten min verademd min productie
$= 0$, alles opgenomen) en niet gegeten productie van de vleeseters
$80 - 20 = \qty{60}{}$; productie van de toproofdieren
\qty{6}{}: samen \qty{7386}{kcal}.
\textbf{15.} $12\,000 + 1100 + 300 + 14 + 7386 = \qty{20800}{kcal}$.
\textbf{16.} Gelijk aan de bruto productie: de bron is in een stationaire
toestand en slaat van jaar tot jaar geen biomassa op.
\textbf{17.} $7386/20\,800 = \qty{36}{\%}$ (het meeste van de rest is de
eigen ademhaling van de \omterm{def:b1:ecosystem-organization:trophic}{producenten}).
\textbf{18.} \omterm{def:b1:ecosystem-organization:trophic}{Producenten} \qty{8800}{}, planteneters \qty{400}{}, vleeseters
\qty{80}{}, toproofdieren \qty{6}{}; rendementen \qty{4.5}{\%},
\qty{20}{\%}, \qty{7.5}{\%}.
\textbf{19.} \qty{8000}{}, \qty{400}{}, \qty{100}{}, \qty{15}{kcal/m^2}:
rechtop, elk niveau een tiende of minder van het niveau eronder.
\textbf{20.} \omterm{def:b1:ecosystem-organization:trophic}{Producenten} $8000/8800 = 0.9$ jaar; planteneters $400/400 =
1$ jaar; vleeseters $100/80 = 1.25$ jaar; toproofdieren $15/6 =
2.5$ jaar. De planten worden het snelst omgezet: kortlevend \omterm{def:b1:body-plans-tissues:tissue}{weefsel} dat
onafgebroken wordt begraasd; de toproofdieren zijn langlevende dieren
waarvan de voorraad traag wordt vernieuwd.
\textbf{21.} Een reiger van \qty{2}{kg} weegt ongeveer \qty{500}{g} droog,
\qty{5000}{kcal}, en heeft ongeveer dat aan productie per jaar nodig om
zichzelf en zijn jongen te vervangen: $5000/6 \approx \qty{800}{m^2}$ bron
--- in de praktijk veel meer, want hij verademt ook: bij
\qty{100}{kcal} per dag is dat \qty{36500}{kcal} opname per jaar, dat wil
zeggen \qty{6000}{m^2} van het niveau eronder (productie van de vleeseters
\qty{80}{}), en zo ziet het territorium van een reiger er ook uit.
\textbf{22.} De extra netto productie wordt niet gegeten en gaat naar de
\omterm{def:b1:ecosystem-organization:trophic}{reducenten}, waarvan de ademhaling verdubbelt; 's nachts, zonder
fotosynthese, kan hun zuurstofvraag het water van zuurstof beroven en de
vissen doden --- eutrofiëring.
\textbf{23.} De productie van de planteneters daalt tot \qty{200}{}: de
vleeseters krijgen de helft van hun voedsel, hun productie zakt naar
\qty{40}{}, die van de toproofdieren naar 3, en er zijn minder reigers; het
zeegras, dat minder wordt begraasd, wordt dichter en gaat voor een groter
deel ongegeten naar de \omterm{def:b1:ecosystem-organization:trophic}{reducenten}.
\textbf{24.} De energie van de toproofdieren is piepklein, maar hun uitwerking
gaat via het aantal vleeseters, dat de planteneters beheerst, die het gras
beheersen: hen wegnemen laat de vleeseters stijgen, de planteneters dalen en
het gras groeien --- een cascade door het web
(\cref{ch:b1:species-interactions}); enkele kilocalorieën planteneters
wegnemen verandert niets dan enkele kilocalorieën.
\textbf{25.} \omterm{def:b1:ecosystem-organization:trophic}{Producenten} naar planteneters \qty{4.5}{\%} (van de NPP;
\qty{16}{\%} van het gegetene), planteneters naar vleeseters \qty{20}{\%},
vleeseters naar toproofdieren \qty{7.5}{\%}; zonlicht naar bruto productie
\qty{1.2}{\%}; enkele miljoensten van het zonlicht eindigt in de
toproofdieren.
\end{solution}
